\documentclass{jsarticle}
\usepackage{amsmath}
\usepackage{amssymb}
\begin{document}
\title{}
\author{}
\date{}
\maketitle

     あけましておめでとうございます。今年もよろしくお願いいたします。


  $$ \bigoplus{(i \hbar^{\nabla})}^{\oplus L} = e^{- x \log x} $$
  上の式は、代数幾何の量子化の式であり、
  この式は、以下の式からの求まりから導かれている。
  $$ H\Psi = \bigoplus{{H\Psi }\over {\nabla L}} $$
  ガンマ関数の大域的積分多様体が、以下の式で表されている。
  $$ = \int \Gamma(\gamma)^{'}dx_m $$
  これらの式は、アカシックレコードの未来の子からもらった式であり、
  母の明王台高等学校時代とアカシックレコードの未来の子からのヒントから、
  大域的微分多様体を求めることが出来たことである。
  $$ {d \over df}F(x) = m(x) $$

    母は梅沢先生でついていて、
  順子お姉さんは、落水壮のライト建築家でついていて、睦世お姉さんは、
  マリー・キュリーでついていて、父は、フェルニー・ハイゼンベルクでついていて、
  匠生君は、アンリ・ポワンカレでついていて、俊咲くんは、ノーベル財団の
  授賞式でついていて、計也さんは、フィールズ財団の創設日でついていて、
  ユウくんは、アルベルトでついていて、
  フルメタをしている家族になっていて、安岡正篤先生は、尊敬の眼差しです。


\end{document}
